\clearpage\thispagestyle{empty}
\pdfbookmark[0]{Sources for Gauss to Bode (3 September 1820)}{bodeseptember1820-sources}
\begin{center}\textbf{Sources for Gauss to Bode (3 September 1820)}\end{center}
\noindent This AI translation follows \textit{Gauss an Bode} in \href{https://gdz.sub.uni-goettingen.de/id/PPN236007467}{\textit{Werke}, vol.~VI} (Göttingen, 1874), pp.~433--434. The \href{https://api.digitale-sammlungen.de/iiif/image/v2/bsb10538352_00234/full/1600,/0/default.jpg}{first printing} appeared under \textit{Astronomische Beobachtungen im Jahr 1820} in the \textit{Astronomisches Jahrbuch für das Jahr 1823}, edited by Johann Elert Bode and issued in Berlin in 1820, pp.~228--230. The links lead to the Bavarian State Library copy. The letter begins below the preceding table. Gauss credits von Staudt with calculating the Pallas ephemeris, which is mentioned in the letter and lies outside the collected text.
\par\smallskip\noindent Gauss refers to the \hyperlink{native-434-part-lower}{report of 5 June 1820} for his description of Reichenbach’s meridian circle. The \href{https://api.digitale-sammlungen.de/iiif/image/v2/bsb10538352_00236/full/1600,/0/default.jpg}{first printing on p.~230} spells the name Bradley. The collected edition spells it Brandley, as in the translation. The tables follow the collected edition’s arrangement and keep the letters O and U for the upper and lower limbs of the Sun. Translation and typesetting produced by AI.
\clearpage
